\begin{abstract}
    Does adding more agents help an LLM workflow once compared systems share the same benchmark loader, tool access, answer contract, usage accounting, and trajectory logging?
    We introduce BenchAgent, an evaluation framework that places single-agent, fixed multi-agent (MAS), and evolving MAS workflows under one normalized execution and logging protocol.
    BenchAgent evaluates these substrate-internal workflows across ten reasoning, coding, and tool-use benchmarks with GPT-4.1, and separately reports a Protocol-Aligned External (PAE) GAIA study of a runtime-generated workflow.
    Under SI conditions, at most one of six tested MAS exceeds the matched single-agent anchor on benchmark-balanced average accuracy---EvoAgent lies within the Wilson one-run guidance---while the remaining five trail by 2.56--11.29 points and occupy more expensive accuracy--cost trade-offs.
    On the PAE GAIA snapshot, a Claude-Code-style runtime workflow reaches 66.72\% overall and 69.23\% on Level~3, more than 20 points above the strongest non-Claude baseline (Jarvis, a fixed MAS).
\end{abstract}


\setlength{\parskip}{1.5pt plus1.5pt minus0pt}


\section{Introduction}
\label{sec:introduction}
LLM-based agent research has expanded from single reasoning-and-acting loops~\citep{yao2022react,schick2023toolformer,nakano2021webgpt,yang2024sweagent} and role-specialized multi-agent systems~\citep{li2023camel,wu2023autogen,hong2023metagpt} to designs that configure, evolve, or generate the workflow itself~\citep{chen2023agentverse,yuan2024evoagent,hu2024adas,zhang2024aflow,fourney2024magenticone}.

Accuracy alone cannot separate workflow organization from protocol advantage.
Recent benchmarks broaden coverage to interactive tasks, tool use, web environments, software engineering, and diagnostic signals~\citep{liu2023agentbench,mialon2023gaia,ma2024agentboard,zhou2023webarena,jimenez2024swebench,qin2023toolbench}.
Yet cross-paradigm comparisons routinely leave inputs, answer contracts, tool surfaces, usage accounting, or trajectory logging uneven across systems.
For example, debate and orchestration frameworks explicitly add extra rounds, role-specialized handoffs, or tool-routing policies relative to a single-controller loop~\citep{du2023multiagentdebate,wu2023autogen,shen2023hugginggpt}; if such systems are compared only by final accuracy, the observed gap can mix coordination with protocol advantage.
Our same-instance contrast in Figure~\ref{fig:gaia_generated_team_failure} later shows a complementary failure mode: an intermediate handoff can drop a task constraint that is invisible in the final answer alone.
This is why analytical traces matter in agent evaluation frameworks~\citep{ma2024agentboard,harbor_framework}: without them, equivalent-looking systems can spend different tokens, lose different constraints, or handle tool failures differently.

\begin{figure*}[!t]
    \centering
    \includegraphics[width=0.92\textwidth]{resources/figure/methodology.drawio.png}
    \vspace{-0.5em}
    \caption{\small
        \textbf{BenchAgent compares workflow paradigms under a shared evaluation substrate}. Benchmark instances enter the same loading, tool-access, accounting, logging, and evaluation interfaces, while the workflow layer varies across single-agent, fixed MAS, dynamic/evolving MAS, and externally evaluated runtime-generated workflows. Within the substrate-internal setting, this design isolates workflow lift from differences in benchmark loading, tool implementation, or evaluation criteria; under the PAE setting, it documents aligned runtime workflows under the same evaluation target. The central question is whether workflow organization itself changes accuracy, cost, and trace behavior.
    }
    \label{fig:paradigm_centered_setup}
\end{figure*}

Figure~\ref{fig:paradigm_centered_setup} sketches the evaluation setup: benchmark instances enter a shared runtime and evaluator, while only the workflow layer varies across single-agent, fixed MAS, dynamic/evolving MAS, and externally evaluated runtime-generated execution.
The goal is not to identify the strongest possible agent controller, but to estimate workflow lift: \emph{whether reorganizing a matched controller into a MAS workflow yields reliable marginal gain under explicit protocol boundaries.}
Accordingly, the paper reports one controlled substrate-internal comparison and one protocol-aligned external runtime case study rather than treating all systems as a single symmetric leaderboard.

We build BenchAgent for this purpose.
BenchAgent is an evaluation framework; it wraps existing benchmarks in a common execution interface and records the process signals needed for workflow comparison.
BenchAgent places single-agent systems, fixed MAS, and evolving MAS under the same benchmark-normalized interface, sandbox runtime, tool-allocation rules, and process-level logger.
We use four workflow labels throughout the paper: single-agent, fixed MAS, dynamic/evolving MAS, and runtime-generated workflow.
We call the first three a \emph{Substrate-Internal} (SI) comparison when the workflow runs inside BenchAgent and exposes its messages, tool calls, usage, and termination state to the same logger.
For systems that cannot be reimplemented without changing their runtime semantics, we use a \emph{Protocol-Aligned External} (PAE) comparison: inputs, evaluator, final-answer format, tool-capability classes, backend model family, and permission profile are aligned, while controller details remain visible through retained traces.

With BenchAgent, we compare a single-agent configuration and representative fixed and evolving MAS with a GPT-4.1 backend.
The broad suite covers reasoning, coding, and tool-use benchmarks (see Table~\ref{tab:main_exp1_overall}).
On GAIA~\citep{mialon2023gaia}, we evaluate the Claude-Code-style runtime workflow (CC-workflow) through a PAE study and compare it with the BenchAgent-based MAS workflows.

The broad benchmark results show that increasing the number of agents or adding explicit coordination does not by itself guarantee positive workflow lift.
Fixed and evolving MAS vary in effectiveness, token usage, and latency, and several trail the matched single-agent anchor.
The GAIA comparison asks a narrower system-level question: whether a mature runtime workflow reaches a different accuracy--cost trade-off on long-horizon tool-use tasks.
It does on Levels~2--3, but this is a deployed-configuration result, not evidence for any single mechanism.
\textbf{Our key contributions are:}
\begin{itemize}[nosep, leftmargin=12pt]
    \item We introduce BenchAgent, a workflow-evaluation substrate that normalizes benchmark loading, tool access, answer contracts, usage accounting, and trajectory logging for cross-paradigm agent comparison.
    \item Under SI conditions with a GPT-4.1 backend, shared tools, evaluator, and logger, at most one of six tested MAS exceeds the matched single-agent anchor on average, and even that case (EvoAgent, +1.44 points) lies within the Wilson one-run guidance; the remaining five trail by 2.56--11.29 points.
    \item We report a PAE GAIA snapshot in which the CC-workflow reaches 66.72\% overall accuracy, leads on Levels~2--3, and uses fewer recorded tokens than the strongest non-Claude baseline; isolating which design choice drives the gain requires controlled ablations beyond our deployed-configuration comparison.
          \looseness=-1
\end{itemize}


\section{Related Work}
\subsection{Agent Workflow Design}
ReAct, Toolformer, WebGPT, SWE-agent, and HuggingGPT/JARVIS show that a single instrumented controller can already serve as a competitive workflow~\citep{yao2022react,schick2023toolformer,nakano2021webgpt,yang2024sweagent,shen2023hugginggpt}.
We therefore treat the single-agent baseline as a matched workflow anchor rather than as a bare model call or as a claim about the strongest possible single-controller design; recent work corroborates that a strong single-agent baseline can match or exceed homogeneous MAS~\citep{xu2026rethinking}.

Multi-agent systems make the division of labor explicit.
CAMEL and MetaGPT define agent roles and communication procedures~\citep{li2023camel,hong2023metagpt}, AutoGen provides a conversable-agent substrate for fixed or adaptive interaction patterns~\citep{wu2023autogen}, and debate-style systems aggregate multiple proposals before finalization~\citep{du2023multiagentdebate,chan2023chateval}.
Li~\emph{et al.}~\citep{li2024moreagents} show that sampling-and-voting improves several LLM benchmarks; DSPy~\citep{khattab2023dspy} instead optimizes LM pipelines through compilation.
Dynamic and evolving systems search over or mutate the workflow structure within a designer-specified space, as in AgentVerse, EvoAgent, ADAS, AFlow, and Magentic-One~\citep{chen2023agentverse,yuan2024evoagent,hu2024adas,zhang2024aflow,fourney2024magenticone}; MaAS extends this by searching over query-dependent agent topologies via an agentic supernet~\citep{zhang2025maas}; MASPO further co-optimizes agent prompts jointly across the pipeline~\citep{wang2026maspo}.
Cemri~\emph{et al.}~\citep{cemri2025maswhyfail} catalog MAS failure modes attributable to coordination design; Kim~\emph{et al.}~\citep{kim2025scaling} further argue that net workflow gain requires coordination benefit to exceed communication overhead.
These systems motivate our taxonomy, yet prior comparisons rarely run single-agent, fixed MAS, evolving MAS, and runtime workflows under a unified logging and accounting protocol.
\looseness=-1

\subsection{Agent Evaluation Frameworks}
AgentBench, GAIA, AgentBoard, WebArena, SWE-bench, ToolBench, and Silo-Bench~\citep{liu2023agentbench,mialon2023gaia,ma2024agentboard,zhou2023webarena,jimenez2024swebench,qin2023toolbench,zhang2026silobench} extend agent evaluation to interaction, tool use, web/software tasks, analytical traces, and distributed coordination.
Harbor~\citep{harbor_framework} is closest in infrastructure, evaluating and optimizing sandboxed agents and models in containerized environments.
BenchAgent occupies a different scope: it introduces neither a new task dataset nor a benchmark split, but instead aligns execution interfaces, tool surfaces, usage accounting, and trajectories, so that coordination policies can be compared on existing benchmarks.
Table~\ref{tab:appendix_framework_feature_comparison} in Appendix summarizes the feature-level distinction from AgentBench and Harbor.

\subsection{Runtime-Generated Workflows}
Recent engineering systems push part of workflow design into the execution loop.
LangGraph~\citep{langgraph_docs} and CrewAI~\citep{crewai_docs} provide graph orchestration and role-based crews; Anthropic~\citep{anthropic_building_effective_agents,anthropic_multiagent_research_system,anthropic_context_engineering} documents subagent creation, context-window isolation, and permission-scoped delegation patterns instantiated in Claude Code~\citep{anthropic_claude_code_product,anthropic_claude_code_overview,anthropic_claude_code_subagents}; similar agentic coding interfaces appear in OpenAI Codex and OpenCode~\citep{openai_codex_product,opencode_product}.
These systems are documented primarily through official documentation rather than peer-reviewed publications; we cite them as engineering descriptions, not empirical evidence.
Our GAIA experiment asks whether the CC-workflow paired with a GPT-4.1 backend reaches a different accuracy--cost profile from BenchAgent-based fixed and evolving MAS under a documented PAE comparison.


\section{Evaluation Protocol}
\label{sec:evaluation_protocol}
\subsection{Workflow Lift as the Measurement Target}
We evaluate agent systems as \emph{workflow organizations}, not isolated products.
For a benchmark instance \(x\), a workflow \(w\) produces a final answer \(\hat{y}=w(x)\) and, when available, an execution trace \(\tau_w(x)\) containing model calls, tool calls, messages, artifacts, and termination events.
A task evaluator \(E(x,\hat{y}) \in \{0,1\}\) measures success, while a cost summary \(c(\tau_w)\) records usage signals such as tokens, latency, tool calls, and delegation structure.
The comparison therefore asks not only whether \(\hat{y}\) is correct, but also what trace \(\tau_w\) produced it.

We distinguish workflow categories by how \(\mathcal{A}_t\) (active agents), \(G_t\) (communication or delegation topology), and \(\mathcal{T}_t\) (tool scopes) are specified over time:
\begin{itemize}[leftmargin=12pt, nosep]
    \item \textbf{Single-agent workflow.} Keeps \(|\mathcal{A}_t|=1\) throughout; a single controller handles the entire trace.
    \item \textbf{Fixed MAS.} Operates with a predefined \((\mathcal{A},G,\mathcal{T})\), such as solver--critic--aggregator or debate-style handoffs.
    \item \textbf{Evolving MAS.} Selects or mutates \(G\) at runtime from a designer-specified family \(\mathcal{G}\) (a.k.a.\ dynamic MAS). \looseness=-1
    \item \textbf{Runtime-generated workflow.} Changes \(\mathcal{A}_t\), \(G_t\), and \(\mathcal{T}_t\) during execution by creating task-specific agents, assigning private context, choosing tool scopes, or adding verification and recovery branches.
\end{itemize}
Table~\ref{tab:appendix_workflow_taxonomy} gives the full terminology map.

\begin{definition}[Workflow lift]
    \label{def:workflow_lift}
    Our primary substrate-internal quantity is \emph{workflow lift}: the change in accuracy and cost when a single-agent workflow is replaced by a fixed or evolving MAS workflow, with the base model, benchmark loader, tool interface, answer contract, evaluator, and accounting substrate held constant.
\end{definition}

This framing does not require the single-agent anchor to be the strongest possible ReAct implementation; instead, it requires the compared workflows to share the same control substrate, so that the measured difference reflects the tested workflow organization. Section~\ref{subsec:pass1_wilson} describes how single-run lift estimates are interpreted under Wilson guidance.


\subsection{Comparison Setup: SI and PAE}
\label{subsec:comparison_setup}
BenchAgent provides the shared execution substrate for the SI comparison and normalizes six interfaces---benchmark loading, input and answer formatting, runtime control, tool access, usage and trajectory logging, and evaluator calls.
In the broad-benchmark experiment (Section~\ref{subsec:broad_bench_results}), Jarvis/HuggingGPT-style execution, LLM-Debate, EvoAgent, AutoGen, CAMEL, and ChatEval are implemented as workflow instances over this substrate rather than as unrelated tool stacks.
This substrate renders final scores, token usage, latency, tool calls, message histories, agent identifiers, and stage-level traces comparable across workflows.
Appendix~\ref{subsec:appendix_prompt_traj_examples} includes retained prompt excerpts that document the controller instructions and output contract used for the external runtime workflow.
\looseness=-1

Claude Code cannot be reimplemented inside BenchAgent without altering its runtime semantics.
The GAIA experiment (Section~\ref{subsec:gaia_results}) therefore adopts a PAE comparison.
The CC-workflow shares the GPT-4.1 backend family, GAIA validation inputs, answer-only output schema, evaluator, and GAIA-relevant tool-capability classes with the SI baselines, while internal controller details remain only partially visible through retained traces.
Table~\ref{tab:appendix_gaia_protocol_summary} in the Appendix lists which protocol fields are aligned and which cannot be certified from the released artifacts.
\looseness=-1

\subsection{Reporting Protocol and Wilson Guidance}
\label{subsec:pass1_wilson}
For each instance, we record task success, token usage, wall-clock latency, and the full trajectory.
When traces expose sufficient structure, we also derive process summaries---spawned agents, tool calls, and delegation depth---used for interpretation, not as additional success metrics.
Broad-benchmark average accuracy is the unweighted mean over ten benchmark-level accuracies; GAIA overall is weighted by level sizes $(N_1,N_2,N_3)=(53,86,26)$.
All numbers are pass@1 single-run results.

We use the \emph{Wilson 95\% binomial confidence interval}~\citep{wilson1927} as a conservative scale for one-run accuracy gaps: for an observed proportion $\hat{p}$ over $N$ instances, the interval covers the unknown success probability more reliably than the normal approximation when $\hat{p}$ is near $0$ or $1$ or when $N$ is small.
We report its two-sided half-width per benchmark in Table~\ref{tab:appendix_wilson_widths} and treat any pass@1 gap smaller than this half-width as descriptive rather than as stable ordering evidence.
Concretely, differences below roughly 17 points on AIME ($N{=}30$) or 18 points on GAIA Level~3 ($N{=}26$) should not be read as ordering evidence under one run.

\section{Experiments}
\label{sec:experiments}
We instantiate two experiments under the protocol of Section~\ref{sec:evaluation_protocol}.
The first experiment (Section~\ref{subsec:broad_bench_results}) tests whether fixed and evolving MAS produce positive workflow lift over a matched single-agent anchor on ten broad benchmarks under SI conditions; the second (Section~\ref{subsec:gaia_results}) places a Claude-Code-style runtime workflow on GAIA under PAE alignment.
Each number reported below is a pass@1 result from one end-to-end run: we treat gaps smaller than the Wilson half-widths in Appendix Table~\ref{tab:appendix_wilson_widths} as descriptive rather than as stable ordering evidence, and we do not attribute the CC-workflow's GAIA advantage to a single mechanism without ablation evidence.

\subsection{Experimental Setup}
\label{subsec:exp_setup}
\noindent\textbf{Datasets.}
The broad SI comparison covers ten benchmarks: MATH, AIME, GSM8K, DROP, BBH, MMLU-Pro, HumanEval, MBPP, HotpotQA, and IFEval.
The runtime-workflow study uses the full GAIA validation split~\citep{mialon2023gaia} ($(N_1,N_2,N_3)=(53,86,26)$, weighted by level size for overall accuracy), which stresses multi-step retrieval, file inspection, state preservation, and delayed final-answer release.

\noindent\textbf{Tool allocation.}
Tool access is benchmark-specific but fixed across compared workflows.
In the SI suite, HotpotQA uses the expanded full BenchAgent tool registry, whereas MATH, AIME, GSM8K, DROP, BBH, MMLU-Pro, HumanEval, MBPP, and IFEval expose only \texttt{python\_interpreter} beyond final-answer emission.
GAIA uses the corresponding full-tool regime: BenchAgent-based systems receive the expanded full registry, and the CC-workflow is matched at the level of GAIA-relevant tool-capability classes under the PAE protocol.
Appendix Table~\ref{tab:appendix_tool_regimes} summarizes the benchmark-level allocation.

\noindent\textbf{Compared systems.}
The single-agent anchor is BenchAgent Core---a full substrate run, not a bare LLM call---sharing the same loader, evaluator, and tool registry as all MAS wrappers.
It serves as the workflow-lift estimator for the broad-benchmark experiment (Section~\ref{subsec:broad_bench_results}), not as a claim about the strongest possible single-agent controller.
Fixed MAS baselines include Jarvis~\citep{shen2023hugginggpt}, LLM-Debate~\citep{du2023multiagentdebate}, AutoGen~\citep{wu2023autogen}, CAMEL~\citep{li2023camel}, and ChatEval~\citep{chan2023chateval}; EvoAgent~\citep{yuan2024evoagent} represents the evolving MAS category.
Implementation notes, transfer sanity checks, and a vanilla ReAct anchor calibration appear in Appendix Tables~\ref{tab:appendix_baseline_fidelity} and~\ref{tab:appendix_react_calibration}.

\noindent\textbf{CC-workflow.}
Evaluated only on GAIA under a PAE protocol: inputs, answer schema, evaluator, backend model family, and GAIA-relevant tool-capability classes are aligned; the controller remains partially external.
Full protocol documentation appears in Appendix Tables~\ref{tab:appendix_gaia_protocol_summary} and~\ref{tab:appendix_gaia_tool_alignment}.

\begin{table*}[!t]
    \centering
    \caption{\small
        \textbf{Substrate-Internal broad-benchmark comparison of single-agent, fixed MAS, and evolving MAS under matched protocols.} Wilson 95\% confidence-interval half-widths are shown for the individual benchmark rows only; the Avg. Acc. row reports benchmark-balanced average accuracy, while Avg. Tok. reports instance-level average end-to-end token usage and Avg. Time reports instance-level average execution time. Section~\ref{subsec:pass1_wilson} describes averaging and interpretation.
    }
    \vspace{-1em}
    \label{tab:main_exp1_overall}
    \setlength{\tabcolsep}{3.8pt}
    \renewcommand{\arraystretch}{1.16}
    \small
    \resizebox{\textwidth}{!}{
        \begin{tabular}{lccccccc}
            \toprule
            \textbf{Benchmark / Metric} & \textbf{Single Agent}          & \textbf{EvoAgent}              & \textbf{LLM-debate}            & \textbf{Camel}                 & \textbf{AutoGen}               & \textbf{Jarvis}                & \textbf{ChatEval}             \\
            \midrule
            MATH                        & 66.75{\scriptsize $\pm$4.6}\%  & 61.75{\scriptsize $\pm$4.7}\%  & 69.34{\scriptsize $\pm$4.5}\%  & 61.05{\scriptsize $\pm$4.8}\%  & 40.53{\scriptsize $\pm$4.8}\%  & 60.25{\scriptsize $\pm$4.8}\%  & 59.75{\scriptsize $\pm$4.8}\% \\
            AIME                        & 46.67{\scriptsize $\pm$16.8}\% & 36.67{\scriptsize $\pm$16.3}\% & 26.67{\scriptsize $\pm$15.1}\% & 26.08{\scriptsize $\pm$15.0}\% & 20.00{\scriptsize $\pm$13.9}\% & 26.70{\scriptsize $\pm$15.1}\% & 3.33{\scriptsize $\pm$8.0}\%  \\
            GSM8K                       & 94.50{\scriptsize $\pm$2.3}\%  & 93.25{\scriptsize $\pm$2.5}\%  & 93.50{\scriptsize $\pm$2.4}\%  & 86.05{\scriptsize $\pm$3.4}\%  & 88.00{\scriptsize $\pm$3.2}\%  & 89.00{\scriptsize $\pm$3.1}\%  & 85.00{\scriptsize $\pm$3.5}\% \\
            DROP                        & 90.50{\scriptsize $\pm$2.9}\%  & 93.00{\scriptsize $\pm$2.5}\%  & 89.75{\scriptsize $\pm$3.0}\%  & 88.62{\scriptsize $\pm$3.1}\%  & 83.25{\scriptsize $\pm$3.7}\%  & 90.00{\scriptsize $\pm$3.0}\%  & 90.50{\scriptsize $\pm$2.9}\% \\
            BBH                         & 78.25{\scriptsize $\pm$4.0}\%  & 94.00{\scriptsize $\pm$2.4}\%  & 71.48{\scriptsize $\pm$4.4}\%  & 61.00{\scriptsize $\pm$4.8}\%  & 66.67{\scriptsize $\pm$4.6}\%  & 53.50{\scriptsize $\pm$4.9}\%  & 73.00{\scriptsize $\pm$4.3}\% \\
            MMLU-Pro                    & 86.50{\scriptsize $\pm$3.4}\%  & 87.00{\scriptsize $\pm$3.3}\%  & 79.50{\scriptsize $\pm$3.9}\%  & 60.11{\scriptsize $\pm$4.8}\%  & 69.50{\scriptsize $\pm$4.5}\%  & 77.50{\scriptsize $\pm$4.1}\%  & 71.00{\scriptsize $\pm$4.4}\% \\
            HumanEval                   & 84.73{\scriptsize $\pm$6.2}\%  & 81.68{\scriptsize $\pm$6.6}\%  & 93.89{\scriptsize $\pm$4.2}\%  & 89.68{\scriptsize $\pm$5.3}\%  & 61.83{\scriptsize $\pm$8.2}\%  & 88.55{\scriptsize $\pm$5.5}\%  & 86.25{\scriptsize $\pm$5.9}\% \\
            MBPP                        & 68.32{\scriptsize $\pm$4.9}\%  & 73.02{\scriptsize $\pm$4.7}\%  & 72.72{\scriptsize $\pm$4.7}\%  & 75.07{\scriptsize $\pm$4.6}\%  & 73.02{\scriptsize $\pm$4.7}\%  & 75.95{\scriptsize $\pm$4.5}\%  & 75.36{\scriptsize $\pm$4.6}\% \\
            HotpotQA                    & 62.25{\scriptsize $\pm$4.7}\%  & 68.50{\scriptsize $\pm$4.5}\%  & 65.00{\scriptsize $\pm$4.7}\%  & 55.53{\scriptsize $\pm$4.8}\%  & 58.75{\scriptsize $\pm$4.8}\%  & 64.75{\scriptsize $\pm$4.7}\%  & 60.00{\scriptsize $\pm$4.8}\% \\
            IFEval                      & 62.75{\scriptsize $\pm$4.7}\%  & 66.75{\scriptsize $\pm$4.6}\%  & 53.79{\scriptsize $\pm$4.9}\%  & 60.50{\scriptsize $\pm$4.8}\%  & 66.73{\scriptsize $\pm$4.6}\%  & 42.75{\scriptsize $\pm$4.8}\%  & 84.25{\scriptsize $\pm$3.6}\% \\
            \midrule
            Avg. Acc.                   & 74.12\%                        & 75.56\%                        & 71.56\%                        & 66.37\%                        & 62.83\%                        & 66.90\%                        & 68.84\%                       \\
            Avg. Tok.                   & 27434.55                       & 34153.68                       & 36669.25                       & 8470.19                        & 13603.05                       & 9668.33                        & 105838.02                     \\
            Avg. Time                   & 106.29s                        & 82.76s                         & 26.63s                         & 15.75s                         & 14.22s                         & 11.20s                         & 56.08s                        \\
            \bottomrule
        \end{tabular}
    }
    \renewcommand{\arraystretch}{1.0}
\end{table*}

\begin{figure*}[!t]
    \centering
    \includegraphics[width=0.92\textwidth]{resources/figure/main_experiment1_overview/tradeoff_accuracy_cost.pdf}
    \vspace{-1em}
    \caption{\small
        \textbf{Accuracy--cost and accuracy--time trade-offs under SI conditions.}
        Left: benchmark-balanced average accuracy against instance-level average end-to-end token usage. Right: benchmark-balanced average accuracy against instance-level average execution time. Each point represents one workflow, and the dashed line traces the empirical Pareto front under this descriptive aggregation.
    }
    \label{fig:main_exp1_cost}
\end{figure*}

\noindent\textbf{Model and infrastructure.}
All reported main-text BenchAgent results use GPT-4.1 (\texttt{gpt-4.1-2025-04-14})~\citep{openai_gpt41_model}.
Dataset sampling, tool-surface summaries, GAIA-only backend-generalization checks, and trajectory evidence appear in the appendix; the backend checks are scoped to Qwen3-32B for BenchAgent-compatible GAIA workflows and GLM-5 for the BenchAgent GAIA anchor (Appendix Tables~\ref{tab:appendix_qwen32b_gaia_generalization} and~\ref{tab:appendix_glm5_gaia_generalization}).

\subsection{Broad Benchmark Results}
\label{subsec:broad_bench_results}
This experiment tests whether fixed or evolving MAS produce workflow lift over the BenchAgent single-agent anchor under matched model, tool surface, evaluator, and logging protocol; results appear in Table~\ref{tab:main_exp1_overall}, with the accuracy--cost view in Figure~\ref{fig:main_exp1_cost}. Appendix Table~\ref{tab:appendix_react_calibration} calibrates the anchor against a standalone vanilla ReAct controller without redefining the comparison target.
\looseness=-1

\begin{table*}[!t]
    \centering
    \caption{\small
        \textbf{GAIA validation pass@1 comparison under a PAE runtime-workflow setting.} Tok. reports instance-level average end-to-end token usage, and Time reports instance-level average wall-clock time.
    }
    \vspace{-1em}
    \label{tab:main_exp2_gaia}
    \setlength{\tabcolsep}{5pt}
    \small
    \renewcommand{\arraystretch}{1.2}
    \resizebox{\textwidth}{!}{
        \begin{tabular}{l l cccccc}
            \toprule
            \textbf{Paradigm}          & \textbf{Method} & \textbf{GAIA-L1}                        & \textbf{GAIA-L2}                       & \textbf{GAIA-L3}                        & \textbf{Avg./Overall}                  & \textbf{Tok.}     & \textbf{Time}    \\
            \midrule
            Single-Agent               & Single Agent    & 58.82{\scriptsize $\pm$12.8}\%          & 30.00{\scriptsize $\pm$9.5}\%          & 19.23{\scriptsize $\pm$14.7}\%          & 37.56{\scriptsize $\pm$7.3}\%          & 459652.83         & 213.33s          \\
            \midrule
            Evolving MAS               & EvoAgent        & 54.90{\scriptsize $\pm$12.9}\%          & 30.00{\scriptsize $\pm$9.5}\%          & 19.23{\scriptsize $\pm$14.7}\%          & 36.30{\scriptsize $\pm$7.3}\%          & 1573918.58        & 325.04s          \\
            Fixed MAS                  & LLM-Debate      & 52.94{\scriptsize $\pm$13.0}\%          & 34.00{\scriptsize $\pm$9.8}\%          & 15.38{\scriptsize $\pm$13.7}\%          & 37.15{\scriptsize $\pm$7.3}\%          & 414812.44         & 201.53s          \\
            Fixed MAS                  & Camel           & 54.90{\scriptsize $\pm$12.9}\%          & 34.00{\scriptsize $\pm$9.8}\%          & 26.92{\scriptsize $\pm$16.2}\%          & 39.60{\scriptsize $\pm$7.4}\%          & 468587.39         & 387.42s          \\
            Fixed MAS                  & AutoGen         & 56.60{\scriptsize $\pm$12.9}\%          & 30.00{\scriptsize $\pm$9.5}\%          & 11.54{\scriptsize $\pm$12.5}\%          & 35.64{\scriptsize $\pm$7.2}\%          & 383190.65         & 188.98s          \\
            Fixed MAS                  & Jarvis          & 66.03{\scriptsize $\pm$12.4}\%          & 43.02{\scriptsize $\pm$10.2}\%         & 19.23{\scriptsize $\pm$14.7}\%          & 46.66{\scriptsize $\pm$7.5}\%          & 332285.96         & 402.11s          \\
            Fixed MAS                  & ChatEval        & 49.02{\scriptsize $\pm$13.0}\%          & 34.88{\scriptsize $\pm$9.9}\%          & 26.92{\scriptsize $\pm$16.2}\%          & 38.17{\scriptsize $\pm$7.3}\%          & 1886517.67        & 443.76s          \\
            \midrule
            Runtime-Generated Workflow & CC-workflow     & \textbf{60.78{\scriptsize $\pm$12.7}\%} & \textbf{69.62{\scriptsize $\pm$9.5}\%} & \textbf{69.23{\scriptsize $\pm$16.7}\%} & \textbf{66.72{\scriptsize $\pm$7.1}\%} & \textbf{52984.69} & \textbf{134.90s} \\
            \bottomrule
        \end{tabular}
    }
    \renewcommand{\arraystretch}{1.0}
\end{table*}

% \looseness=-1
\noindent\textbf{Fixed and evolving MAS do not consistently produce positive workflow lift.} The matched single-agent anchor reaches 74.12\% benchmark-balanced average accuracy, while EvoAgent reaches 75.56\%. These broad-suite averages are descriptive benchmark-balanced means rather than interval estimates on a single shared instance pool; per-benchmark Wilson half-widths in Table~\ref{tab:main_exp1_overall} carry the row-level uncertainty. EvoAgent is the only MAS with numerically positive average lift, but its +1.44-point gain is smaller than the one-run uncertainty guidance; the remaining MAS fall below the anchor at 62.83\%--71.56\%. Holding model, tools, evaluator, and logger fixed, adding roles or handoffs did not lift overall performance. \emph{\textbf{Takeaway:} at most one MAS exceeds the anchor on average, and that case sits within the one-run uncertainty scale.}

\noindent\textbf{Cost varies more sharply than accuracy.} Figure~\ref{fig:main_exp1_cost} shows workflows with similar benchmark-balanced accuracy but different instance-level token usage and latency. EvoAgent's +1.44-point gain comes at higher token cost; ChatEval is far more token-intensive while trailing the baseline; Camel and Jarvis are lighter and faster but lose accuracy. \emph{\textbf{Takeaway:} accuracy-similar workflows can have sharply different token use and latency.}

% \vspace{-0.35em}
\noindent\textbf{Workflow structure, not agent count, explains MAS variation} (Appendix Table~\ref{tab:appendix_process_diagnostics}). Agent count does not explain MAS behavior in our suite: EvoAgent gains on BBH through prompt-scaffold search, LLM-Debate gains on HumanEval and MATH where proposals are verifiable, and ChatEval gains on IFEval under multi-judge instruction checking. The same protocols hurt elsewhere. \emph{\textbf{Takeaway:} aggregate MAS behavior follows task-protocol fit rather than agent count.}

% \vspace{-0.35em}
\noindent\textbf{MAS gains are task-dependent.} EvoAgent's BBH score (94.00\% [$\pm$2.4] vs.\ 78.25\% [$\pm$4.0] for single-agent) reflects task-specific scaffold selection, not a general multi-agent advantage, and raises a fairness question about prompt-search budget equity. LLM-Debate leads on HumanEval and MATH, where proposals are easy to check; ChatEval scores highest on IFEval but collapses on AIME. MAS help when their protocol matches the task error mode, and add overhead or lose information otherwise. Additional task-level views appear in Appendix Figures~\ref{fig:appendix_accuracy_heatmap}--\ref{fig:appendix_delta_lollipops}. \emph{\textbf{Takeaway:} MAS protocols help when they match a task's error mode, and hurt otherwise.}

\subsection{GAIA Runtime Workflow Results}
\label{subsec:gaia_results}
This experiment focuses on GAIA, where long-horizon tool use tests state preservation, evidence management, and recovery after partial failures; results appear in Table~\ref{tab:main_exp2_gaia}, with backend-sensitivity checks in Appendix Tables~\ref{tab:appendix_qwen32b_gaia_generalization} and~\ref{tab:appendix_glm5_gaia_generalization} and the same-instance contrast in Figure~\ref{fig:gaia_generated_team_failure}. The comparison remains PAE, as described in Section~\ref{subsec:comparison_setup} and Appendix Table~\ref{tab:appendix_gaia_protocol_summary}, so it compares aligned deployed configurations rather than isolating an internal mechanism.

\begin{figure*}[t]
    \centering
    \includegraphics[width=0.88\textwidth]{resources/figure/evo_error_claude_true.pdf}
    \vspace{-0.8em}
    \caption{\small
        \textbf{Same-instance GAIA contrast: EvoAgent vs.\ CC-workflow.} EvoAgent loses a task-critical constraint during linear decomposition, whereas the retained CC-workflow trace preserves intermediate state and verifies before finalization.
    }
    \label{fig:gaia_generated_team_failure}
\end{figure*}

\noindent\textbf{The CC-workflow advantage scales with task length.} The CC-workflow reaches 66.72\% ($\pm$7.1) overall, 20.06 points above the strongest non-Claude result (Jarvis, 46.66\%). Jarvis leads on Level~1, where short retrieval chains often suffice, but the CC-workflow leads by 26.60 points on Level~2 and 42.31 points on Level~3, gaps larger than the uncertainty scale in Appendix Table~\ref{tab:appendix_wilson_widths}. Retained accounting also reports fewer tokens and less wall-clock time than the strongest non-Claude baseline, though cost summaries depend on log visibility and provider-side bookkeeping. \emph{\textbf{Takeaway:} the CC-workflow leads by 20.06 points overall, with a 42.31-point Level~3 gap.}

\noindent\textbf{Retained traces expose structured runtime coordination.} The CC-workflow traces show task-specific subagents, reusable evidence artifacts, and verifier-stage events before final-answer release (6.25 subagents on average, max delegation depth 2.5; Appendix Table~\ref{tab:appendix_claude_trj_stats}). Figure~\ref{fig:gaia_generated_team_failure} illustrates the contrast: EvoAgent loses a task-critical constraint and forces an answer, while the CC-workflow preserves state and verifies before finalization. \emph{\textbf{Takeaway:} subagent count and delegation depth (6.25 / 2.5 max) track task structure.}
Appendix~\ref{subsec:appendix_prompt_traj_examples} gives a retained event-level excerpt from the GAIA trajectory logs.

\noindent\textbf{The two failure modes are structurally different.} In fixed and evolving MAS baselines, a failed sub-step typically stays in the same linear handoff: a solver makes a claim, an aggregator receives a compressed version, and the final answer can be forced with a missing constraint. In retained CC-workflow traces, failures appear as local tool or state problems addressable by changing tool strategy, re-reading an artifact, or invoking a verifier. Candidate mechanisms---runtime delegation, persistent artifacts, verifier-stage control, and context management---remain confounded; Appendix Table~\ref{tab:gaia_mechanism_evidence} lists the needed ablations. \emph{\textbf{Takeaway:} baseline failures get baked into linear handoffs; CC-workflow failures stay local and recoverable.}

\section{Discussion}
The broad-benchmark results show task-specific failure modes, not a uniform MAS advantage. Debate-like aggregation helps when independent proposals are checkable (LLM-Debate on HumanEval and MATH), while evolutionary search helps when prompt or role variation finds a better reasoning mode (EvoAgent on BBH). The same machinery hurts when tasks require exact instruction following or tight evidence control, because handoffs can compress context and hide constraints. Table~\ref{tab:main_exp1_overall} is therefore a workflow-lift comparison around a shared SI anchor, not a leaderboard over all single-agent designs.

Runtime-generated workflows require a different evaluation target. The CC-workflow exposes candidate mechanisms---separate context windows, permission scopes, persistent artifacts, and runtime delegation---that remain confounded with provider-side bookkeeping. Evaluations should therefore report workflow topology, context boundaries, tool scopes, and repair gates alongside final scores; otherwise, similar accuracy can hide incompatible coordination policies, costs, and failure modes.

We keep SI and PAE separate: SI isolates workflow wrappers over a common core, while PAE reports what a mature runtime workflow achieves under a documented but less internally controllable protocol. Collapsing them into one leaderboard would mix workflow lift with protocol advantage.

\section{Conclusion}
BenchAgent evaluates LLM-agent workflows under aligned execution and instrumentation. Across broad benchmarks, fixed and evolving MAS do not consistently beat the matched single-agent anchor and occupy different accuracy--cost trade-offs. On GAIA, the CC-workflow performs better on harder levels with lower retained-token and wall-clock cost under PAE, but this remains a deployed-configuration comparison rather than a mechanism estimate. We release BenchAgent and trajectory tooling for joint study of workflow generation, context management, tool scoping, and verification as properties distinct from agent count.

\section{Limitations}
Our comparison is controlled but not fully causal. The broad-benchmark MAS are re-instantiations, not exact reproductions; the CC-workflow is an external engineering exemplar routed to the same backend, so its advantage may combine workflow generation, context compaction, file and shell tooling, permission scoping, and provider-side bookkeeping. The results compare deployed configurations, not isolated mechanisms.

Mechanism evidence is partial: we report one-run pass@1 results and rely on retained traces that expose only part of Claude Code's internals. Appendix Table~\ref{tab:appendix_react_calibration} adds a vanilla ReAct calibration on MATH, BBH, HumanEval, and HotpotQA to contextualize the matched anchor. Future work should add repeated trials, controller-sensitivity checks, stricter tool-surface-matched ablations, runtime-workflow ablations, and fuller instrumentation for context packaging and recovery behavior.
